\addtocontents{toc}{\protect\Needspace{3\baselineskip}}
\originalchapter{光锥量子化与玻色弦谱}\label{ch:spectrum}
\originalsection{光锥坐标与物理模式}
定义 $X^\pm=(X^0\pm X^{D-1})/\sqrt2$, 度规为 $\dd s^2=-2\dd X^+\dd X^-+\dd X^I\dd X^I$, $I=1,\ldots,D-2$. $X^+$ 作世界面时钟, 选开弦 $X^+=x^++2\alpha'p^+\tau$, $p^+\ne0$. 由约束解出 $X^-$ 的导数, 独立振子只剩横向 $D-2$ 个. 这使负范数时向振子不进入物理态空间, 但量子 Lorentz 一致性仍须检查.

Neumann 开弦的正规模展开为
\begin{equation}\label{eq:openmodes}
X^I=x^I+2\alpha'p^I\tau+
\ii\sqrt{2\alpha'}\sum_{n\ne0}\frac{\alpha_n^I}{n}\ee^{-\ii n\tau}\cos n\sigma,
\quad0\le\sigma\le\pi.
\end{equation}
实条件 $\alpha_{-n}=(\alpha_n)^\dagger$. 共轭动量 $\Pi^I=\dot X^I/(2\pi\alpha')$, 等时正则对易关系给 $[x^I,p^J]=\ii\delta^{IJ}$、$[\alpha_m^I,\alpha_n^J]=m\delta_{m+n,0}\delta^{IJ}$. $n>0$ 的 $\alpha_n$ 是湮灭算符, $\alpha_{-n}$ 是产生算符, 归一化振子为 $\alpha_n/\sqrt n$.

\originalsection{正规序与临界维数}
横向 Virasoro 算符 $L_n=\frac12\sum_m:\alpha_{n-m}^I\alpha_m^I:$ 满足 $[L_m,L_n]=(m-n)L_{m+n}+c(m^3-m)\delta_{m+n,0}/12$, $c=D-2$. 中心项来自把无限振子乘积正规序, 而不是经典 Poisson 括号. 逐个模式的零点能为 $(D-2)\sum_{n>0}n/2$; 使用与共形对称相容的解析正则化 $\sum n=-1/12$ 给截距 $a=(D-2)/24$.

另一种共变量子化包括 $D$ 个坐标场和重参数化鬼场, 后者中心荷为 $-26$. Weyl 异常消失要求总中心荷零, 从而玻色弦 $D=26$, 这时光锥截距 $a=1$. 鬼场中心荷及完整 BRST 物理态构造作为进一步工具引用; 本卷在入选作业中证明横向 Virasoro 代数, 不把一个正规化求和独自当成临界维数的完整证明.

\originalsection{开弦与闭弦}
开弦质量壳为 $\alpha'M^2=N-a$, $N=\sum_{n>0}\alpha_{-n}^I\alpha_n^I$. 在 $D=26$, 基态 $M^2=-1/\alpha'$ 是 tachyon, 表示所展开的真空不稳定. 它不能被解释成稳定物体可携信息超光速. $N=1$ 的 24 个横向态 $\alpha_{-1}^I|0;k\rangle$ 对应无质量矢量. 高能级振子分拆计数由 $\prod_{n\ge1}(1-q^n)^{-24}$ 给出.

闭弦有独立左右振子 $\alpha_n^I,\tilde\alpha_n^I$, 周期性要求 $N=\tilde N$. 非紧致自由谱 $\alpha'M^2=4(N-1)=4(\tilde N-1)$. 无质量层 $\alpha_{-1}^I\tilde\alpha_{-1}^J|0\rangle$ 分解为对称无迹、反对称及迹: 分别是引力子、二形式场和 dilaton. 态空间出现自旋 2 是把弦与引力联系起来的入口, 其非线性背景动力学还需世界面一致性与相互作用分析.

\originalsection{光锥能量与模式归一化}
模展开的 $\ii/n$ 不能随意删掉. 对式\eqref{eq:openmodes}求 $\tau$ 导数, 振子部分变为 $\sqrt{2\alpha'}\sum\alpha_n\ee^{-\ii n\tau}\cos n\sigma$. 与 $\Pi=\dot X/(2\pi\alpha')$ 配对时, 正弦或余弦的正交积分提取模式, 正则对易关系正好产生 $[\alpha_m,\alpha_n]=m\delta_{m+n,0}$. $1/n$ 控制量纲与导数, $\ii$ 使正负模式对互为厄米共轭后坐标为实.

零模的质量也可无需猜测截距位置重建. $L_0=\alpha'p_I^2+N$, 光锥约束 $\alpha_0^-=(L_0-a)/(\sqrt{2\alpha'}p^+)$ 且 $\alpha_0^-=\sqrt{2\alpha'}p^-$, 所以 $2p^+p^-=p_I^2+(N-a)/\alpha'$. 由目标时空不变量 $M^2=2p^+p^--p_I^2$ 立即得到谱. 当中心项、截距或模式规范化改变时, 必须把这三步一起重新核查, 不能只调一个质量公式以匹配记忆中的答案.
